\section{Conclusion}
\label{sec:conclusion}

On sparse booked bank accounts, the empirical distribution of an account's
own past balance changes is a hard baseline for a state-space forecaster.
\begin{itemize}
\item \textbf{The state-space model wins at a week.} Its intervals are
  the more cautious about the downside.
\item \textbf{The empirical forecast wins at a month and beyond.}
\item \textbf{Why is not established (exploratory).} Replacing past
  scheduled flows with the projection raises CRPS at a month and beyond,
  but even the realized future flows do not remove the loss at three
  months, and within-window dependence explains only a small part. A
  forward-period confirmatory test on held-out accounts confirmed both,
  and that the replacement does relatively better where schedules are
  predictable.
\item \textbf{New crossings of zero are over-predicted by all four
  model-based forecasts,} and none beats the account's own historical
  frequency at any horizon. Short-horizon negative balances are mostly
  existing ones.
\item \textbf{The planned calibration and sharpness claims fail in
  part,} and the paper says where (Table~\ref{tab:claims}).
\end{itemize}
Two directions follow. One is a hybrid that uses the schedule only when
it has been regular, and the account's empirical changes otherwise,
selected on calibration at every level. The other is a test on modern,
card-heavy accounts with pending transactions, which these data cannot
provide.

\begin{table}[t]
\centering\footnotesize
\begin{tabular}{p{0.07\linewidth}p{0.40\linewidth}p{0.42\linewidth}}
\toprule
Claim & Test fixed in the plan & Outcome \\
\midrule
A1 & Synthetic, model-matched: every interval within $\pm\R{c/a1-tol-off}$
  points (detector off) and $\pm\R{c/a1-tol-on}$ (on) & \textbf{Holds} in
  every cell (Table~\ref{tab:a1}) \\
A2 & Held-out: every interval within $\pm\R{c/a2-tol}$ points at \hA\ and
  \hB\ days, $\pm\R{c/a2-tol-90}$ at \hC & \textbf{Fails}:
  \R{a2/count/fail} cells fail (\lA\% at every horizon, \lB\% at \hB\
  days), \R{a2/count/pass} pass, \R{a2/count/inconclusive} inconclusive
  (Table~\ref{tab:a2}) \\
B1 & Held-out: CRPS skill $>0$ against each baseline at \hB\ and \hC\
  days & \textbf{Holds} against B-SA+, B-30, B-LIN, B-SA, B-SEAS;
  \textbf{fails} against B-EQ and B-GLOBAL, under either aggregation \\
B2 & Realistic generator: as B1 & \textbf{Holds} against every
  per-account baseline; \textbf{fails} against B-GLOBAL (no significant
  difference) \\
C1 & Byte-identical reports on every target & \textbf{Falsified} for the
  frozen engine. After a post-freeze fix, holds on an arm64 Mac (native
  debug and release, UniFFI, WASM under Node) and on x86\_64 Linux (native,
  WASM); not shown on a phone or with other compilers \\
D1 & Predicted sign of each ablation & Holds for lumps, spend floor and
  adaptive noise (coverage). Reversed for the detector, which by the
  plan's rule falsifies it; the alarm-defined strata make this descriptive
  (\S\ref{sec:regime}) \\
\bottomrule
\end{tabular}
\caption{The plan's claims and their outcomes. The primary model for A2
and B1 is V2; for A1, B2 and D1 it is V0, which on synthetic data equals
V2.}
\label{tab:claims}
\end{table}
